\chapter{Лабораторная работа №3. Создание коммутируемой сети Ethernet.}
\label{ch:chapter 3}

\section{\textit{Лабораторная работа №3.1.Основы Ethernet и конфигурирование VLAN}}

\subsection{Цели}

Лабораторная работа помогает получить практические навыки по изучению
следующих тем:

\begin{itemize}

\item Создание VLAN

\item Конфигурирование портов доступа, магистральных портов и гибридных портов

\item Конфигурирование VLAN на основе портов

\item Конфигурирование VLAN на основе MAC-адресов

\item Просмотр таблицы MAC-адресов и информации о VLAN

\end{itemize}

\subsection{Топология сети}

Компании необходимо разделить сеть уровня 2 на несколько VLAN для
удовлетворения служебных требований. Кроме того, VLAN 10 должна обеспечивать
более высокий уровень безопасности, поэтому в нее можно добавить только
специальные ПК.
Для этого пользовательские порты идентичных служб на S1 и S2 необходимо
назначить в одну и ту же VLAN, а порты с определенными MAC-адресами на S2 — в
другую VLAN.

\includegraphics[width=0.5\textwidth]{lb3_1.png} 
\captionof{figure}{Топология сети для лабораторной работы №3}

\subsection{План работы}

\begin{itemize}

\item 1. Создание VLAN.

\item 2. Конфигурирование VLAN на основе портов.

\item 3. Конфигурирование VLAN на основе MAC-адресов.

\end{itemize}

\subsection{Процедура конфигурирования}

\textbf{Шаг 1.} Настроим имена S1 и S2 и отключим ненужные порты.

Зададим имена:

\includegraphics[width=0.5\textwidth]{lb3_2.png} 

\includegraphics[width=0.5\textwidth]{lb3_3.png} 

Отключим порты GE0/0/11 и GE0/0/12 на S1:

\includegraphics[width=0.5\textwidth]{lb3_4.png} 

\textbf{Шаг 2.} Настроим IP-адреса.

Установим для R1 и R3 IP-адреса 10.1.2.1/24 и 10.1.10.1/24 соответственно:

\includegraphics[width=0.5\textwidth]{lb3_5.png} 

\includegraphics[width=0.5\textwidth]{lb3_6.png} 

Установим для S3 и S4 IP-адреса 10.1.3.1/24 и 10.1.3.2/24 соответственно:

\includegraphics[width=0.5\textwidth]{lb3_7.png} 

В данном случае интерфейсы коммутаторов не поддерживают переключение из 2 уровня в 3й, поэтому нужно установить для VLANIF3 на S3 и S4 IP-адреса:

Создадим VLAN 3 на S3 и S4:

\includegraphics[width=0.5\textwidth]{lb3_8.png} 

Настроим порты на S3 и S4 в качестве портов доступа и назначьте их в
соответствующие VLAN:

\includegraphics[width=0.5\textwidth]{lb3_9.png} 

\includegraphics[width=0.5\textwidth]{lb3_10.png} 

Создадим интерфейсы VLANIF и настройте IP-адреса:

\includegraphics[width=0.5\textwidth]{lb3_11.png} 

С помощью команды interface vlanif vlan-id можно создать интерфейс VLANIF и
перейти в режим конфигурирования интерфейса VLANIF.

\includegraphics[width=0.5\textwidth]{lb3_12.png} 

\textbf{Шаг 3.} Создадим VLAN.

Создадим VLAN 2, 3 и 10 на S1 и S2:

\includegraphics[width=0.5\textwidth]{lb3_13.png} 

Команда \textbf{vlan} позволяет создать VLAN, в свою очередь \textbf{batch} позволяет указать для создания несколько VLAN.

\includegraphics[width=0.5\textwidth]{lb3_14.png} 

\textbf{Шаг 4.} Настроим сети VLAN на основе портов.

Настроим пользовательские порты на S3 и S4 в качестве портов доступа и
назначим их в соответствующие VLAN:

\includegraphics[width=0.5\textwidth]{lb3_15.png} 

С помощью команды port link-type { access | hybrid | trunk } можно задать тип
интерфейса, который может быть Access, Trunk или Hybrid.

Команда port default vlan vlan-id позволяет настроить VLAN по умолчанию для
интерфейса и назначить интерфейс в эту VLAN.

\includegraphics[width=0.5\textwidth]{lb3_16.png} 

Настроим порты, соединяющие S1 и S2, в качестве магистральных портов и
разрешите прохождение только пакетов из VLAN 2 и VLAN 3

\includegraphics[width=0.5\textwidth]{lb3_17.png} 

Команда port trunk allow-pass vlan позволяет назначить магистральный порт в
определенные сети VLAN.

\includegraphics[width=0.5\textwidth]{lb3_18.png} 

Команда undo port trunk allow-pass vlan позволяет удалить магистральный порт из
определенных сетей VLAN.

По умолчанию VLAN 1 находится в списке разрешенных сетей. Если VLAN 1 не
используется какой-либо службой, ее необходимо удалить в целях безопасности.

\includegraphics[width=0.5\textwidth]{lb3_19.png} 

\textbf{Шаг 5.} Конфигурация сети VLAN на основе MAC-адресов.

Как показано на схеме сети, R3 имитирует специальный служебный ПК. Допустим, что
данный ПК имеет MAC-адрес a008-6fe1-0c46. Предполагается, что ПК будет
подключаться к сети через любой из портов GigabitEthernet0/0/1, GigabitEthernet0/0 /2
и GigabitEthernet0/0/3 на S2 и передавать данные через VLAN 10

Настроим на S2 привязку MAC-адреса ПК к VLAN 10

Принадлежность к VLAN зависит от исходных MAC-адресов пакетов, и
соответственно добавляются теги VLAN. Этот метод назначения VLAN не зависит от
местоположения, обеспечивая более высокий уровень безопасности и гибкости.

\includegraphics[width=0.5\textwidth]{lb3_20.png} 

Команда mac-vlan mac-address позволяет установить привязку MAC-адреса к VLAN.

Настроим GigabitEthernet0/0/1, GigabitEthernet0/0/2 и GigabitEthernet0/0/3 на S2 в
качестве гибридных портов и разрешите прохождение пакетов из VLAN на основе
MAC-адресов:

На портах доступа и магистральных портах назначение VLAN на основе MAC-адресов
можно использовать только в том случае, если VLAN соответствует PVID. Поэтому
рекомендуется настроить назначение VLAN на основе MAC-адресов на гибридном
порте для получения из нескольких VLAN нетегированных пакетов.

\includegraphics[width=0.5\textwidth]{lb3_21.png} 

Команда port hybrid untagged vlan позволяет назначить гибридный порт в
определенные сети VLAN, чтобы передавать нетегированные кадры.

\includegraphics[width=0.5\textwidth]{lb3_22.png} 

Настроим на портах, соединяющих S1 и S2, разрешение на прохождение пакетов из
VLAN 10

Порты должны разрешать прохождение тегированных кадров из нескольких VLAN.
Следовательно, порты можно настроить в качестве магистральных портов.

\includegraphics[width=0.5\textwidth]{lb3_23.png} 

Настроим S2 и включим назначение VLAN на основе MAC-адресов на GE0/0/1,
GE0/0/2 и GE0/0/3.

Чтобы включить на порте передачу пакетов на основе привязки между MAC-адресом
и VLAN, необходимо выполнить команду \textbf{mac-vlan enable.}.

\includegraphics[width=0.5\textwidth]{lb3_24.png} 

Команда mac-vlan enable позволяет включить функцию назначения VLAN на основе
MAC-адреса для порта.

\includegraphics[width=0.5\textwidth]{lb3_25.png} 

\textbf{Шаг 6.} Выведем на экран информацию о конфигурации.

Выведем на экран информацию о VLAN на коммутаторе:

\includegraphics[width=0.5\textwidth]{lb3_26.png} 

Команда display vlan позволяет вывести на экран информацию о сетях VLAN.

С помощью команды display vlan verbose можно вывести на экран подробную
информацию определенной VLAN, включая идентификатор, тип, описание и состояние VLAN, состояние функции статистики трафика, порты VLAN и режим, в
котором осуществляется назначение портов в VLAN.

\includegraphics[width=0.5\textwidth]{lb3_27.png} 

Выведем на экран конфигурацию назначения VLAN на основе MAC-адресов,
имеющуюся на коммутаторе:

\includegraphics[width=0.5\textwidth]{lb3_28.png} 

Команда display mac-vlan позволяет вывести на экран конфигурацию назначения
VLAN на основе MAC-адресов.

\subsection{Вопросы.} 

Как показано на следующем рисунке, для обеспечения информационной
безопасности определенной услуги только специальные ПК могут получить
доступ к сети через VLAN 10 Как это требование можно реализовать на S1?

Через MAC-vlan

\section{\textit{Лабораторная работа №3.2. Протокол связующего дерева (STP)}}

\subsection{Цели}

Лабораторная работа помогает получить практические навыки по изучению
следующих тем:

\begin{itemize}

\item Включение и отключение STP/RSTP.

\item Процедура изменения режима STP коммутатора.

\item Процедура изменения приоритетов мостов для управления выбором корневого
моста.

\item Процедура изменения приоритетов портов для управления выбором корневого
порта и назначенного порта.

\item Процедура изменения стоимости портов для управления выбором корневого
порта и назначенного порта.

\item Процедура настройки граничных портов.

\item Включение и отключение RSTP.

\end{itemize}

\subsection{Топология сети}

В целях повышения сетевой доступности, компании необходимо создать резервные
каналы в своей коммутируемой сети уровня 2 Кроме того, необходимо развернуть
протокол STP, чтобы предотвратить образование петель на резервных каналах,
вызывающих широковещательный шторм и нестабильность MAC-адресов.

\includegraphics[width=0.5\textwidth]{lb3_29.png} 
\captionof{figure}{Топология сети для лабораторной работы №3}

\subsection{План работы}

\begin{itemize}

\item 1. Включение STP.

\item 2. Изменение приоритетов мостов, чтобы контролировать выбор корневого моста.

\item 3. Изменение параметров порта, чтобы определить роль порта.

\item 4. Изменение протокола на протокол RSTP.

\item 5. Настройка граничных портов.

\end{itemize}

\subsection{Процедура конфигурирования}

\textbf{Шаг 1.} Отключим ненужные порты.

Отключим порт GigabitEthernet0/0/12 между S1 и S2:

\includegraphics[width=0.5\textwidth]{lb3_30.png} 

\includegraphics[width=0.5\textwidth]{lb3_31.png} 

\textbf{Шаг 2.} Включение STP.

Включим STP глобально с помощью команды \textbf{stp enable}, позволяющей включить протокол STP, RSTP или MSTP на коммутационном устройстве или порте:

\includegraphics[width=0.5\textwidth]{lb3_32.png} 

Изменим режим связующего дерева на STP:

\includegraphics[width=0.5\textwidth]{lb3_33.png} 

С помощью команды stp mode{mstp | rstp | stp} можно установить режим работы
протокола связующего дерева на коммутационном устройстве. По умолчанию
коммутационное устройство работает в режиме MSTP. Режим связующего дерева
текущего устройства был изменен на STP.

\includegraphics[width=0.5\textwidth]{lb3_34.png} 

\includegraphics[width=0.5\textwidth]{lb3_35.png} 

\includegraphics[width=0.5\textwidth]{lb3_36.png} 

Выведем на экран статус связующего дерева для S1:

\includegraphics[width=0.5\textwidth]{lb3_37.png} 

Выведем краткую информацию о STP, уже на всех коммутаторах:

\includegraphics[width=0.5\textwidth]{lb3_38.png} 

\includegraphics[width=0.5\textwidth]{lb3_39.png} 

\includegraphics[width=0.5\textwidth]{lb3_40.png} 

\includegraphics[width=0.5\textwidth]{lb3_41.png} 

На основании идентификатора корневого моста и информации о порте каждого
коммутатора текущая топология выглядит следующим образом:

\includegraphics[width=0.5\textwidth]{lb3_42.png} 

\textbf{Шаг 3.} Изменим параметры устройства, чтобы сделать S1 корневым мостом, а S2 —
резервным корневым мостом.

Изменим приоритеты мостов S1 и S2:

\includegraphics[width=0.5\textwidth]{lb3_43.png} 

\includegraphics[width=0.5\textwidth]{lb3_44.png} 

Так как корневой мост играет очень важную роль в сети, то в качестве него обычно
выбирается коммутатор с высокой производительностью и высоким уровнем сетевой
иерархии. Однако такое устройство может иметь невысокий приоритет. Поэтому
необходимо установить коммутатору высокий приоритет, чтобы он мог быть выбран в
качестве корневого моста. С помощью команды stp root можно настроить коммутатор
в качестве корневого моста или резервного корневого моста связующего дерева.

Команда stp root primary позволяет задать коммутатор в качестве корневого
коммутационного устройства. В этом случае коммутатор получит приоритет в
связующем дереве, равный 0, и его нельзя будет изменить.

Команда stp root secondary позволяет задать коммутатор в качестве резервного
корневого моста. В этом случае коммутатор получит приоритет, равный 4096, и
его нельзя будет изменить.

Выведем стату STP для S1:

\includegraphics[width=0.5\textwidth]{lb3_45.png} 

Выведем краткую информацию во всех устройствах:

\includegraphics[width=0.5\textwidth]{lb3_46.png} 

\includegraphics[width=0.5\textwidth]{lb3_47.png} 

\includegraphics[width=0.5\textwidth]{lb3_48.png} 

\includegraphics[width=0.5\textwidth]{lb3_49.png} 

На основании идентификатора корневого моста и информации о порте каждого
коммутатора текущая топология выглядит следующим образом:

\includegraphics[width=0.5\textwidth]{lb3_50.png} 

\textbf{Шаг 4.} Изменим параметры устройства, чтобы назначить порт GigabitEthernet0/0/2
коммутатора S4 корневым портом.

Выведем на экран информацию STP на S4:

\includegraphics[width=0.5\textwidth]{lb3_51.png} 

Изменим стоимость STP порта GigabitEthernet 0/0/1 коммутатора S4 на 50000:

\includegraphics[width=0.5\textwidth]{lb3_52.png} 

Выведем на экран краткую информацию о статусе STP:

\includegraphics[width=0.5\textwidth]{lb3_53.png} 

Выведем на экран информацию о текущем статусе STP:

\includegraphics[width=0.5\textwidth]{lb3_54.png} 

Текущая топология выглядит следующим образом:


\includegraphics[width=0.5\textwidth]{lb3_55.png} 

\textbf{Шаг 5.} Изменим режим связующего дерева на RSTP.

Изменим режим связующего дерева на всех устройствах:

\includegraphics[width=0.5\textwidth]{lb3_56.png} 

\includegraphics[width=0.5\textwidth]{lb3_57.png} 

\includegraphics[width=0.5\textwidth]{lb3_58.png} 

\includegraphics[width=0.5\textwidth]{lb3_59.png} 

Выведем на экран статус связующего дерева. В данном случае для примера
используется S1:

\includegraphics[width=0.5\textwidth]{lb3_60.png} 

\textbf{Шаг 6.} Настроим граничные порты.

Порты GigabitEthernet 0/0/10-0/0/24 коммутатора S3 подключены только к
терминалам, поэтому их необходимо настроить в качестве граничных портов:

Устройство предоставляет несколько портов Ethernet, многие из которых имеют
одинаковую конфигурацию. Настраивать их по очереди утомительно и чревато
возможными ошибками. Самый простой способ — добавить эти порты в группу портов
и выполнить настройку всей группы сразу. Система автоматически выполнит команды
на всех портах в группе.

Команда stp edged-port enable позволяет задать текущий порт в качестве граничного
порта. Если после настройки граничный порт коммутационного устройства получает
BPDU, коммутационное устройство автоматически отменяет настройки порта в
качестве граничного порта и пересчитывает связующее дерево.

В данном случае данная функция недоступна.

\textbf{Вопросы.}

1. Можно ли будет достичь желаемого результата, если на шаге 3 изменить
стоимость GigabitEthernet 0/0/14 коммутатора S1 на 50000? Почему?

-Нет, т.к. мы ставим очень низкий приоритет, из-за чего попросту отключится этот порт для передачи.

2. Как необходимо изменить конфигурацию в текущей топологии, чтобы назначить
порт GigabitEthernet0/0/11 коммутатора S2 корневым портом?

-Изменить стоимость порта GigabitEthernet0/0/10, чтобы отдать приоритет GigabitEthernet0/0/11.

3. Могут ли два канала между S1 и S2 одновременно находиться в режиме передачи
данных? Почему?

-Нет, т.к. тогда может случиться широковещательный шторм.

\section{\textit{Лабораторная работа №3.3. Агрегирование каналов Ethernet.}}

\subsection{Цели}

Лабораторная работа помогает получить практические навыки по изучению
следующих тем:

\begin{itemize}

\item Ручная настройка агрегирования каналов.

\item Настройка агрегирования каналов в статическом режиме LACP.

\item Определение активных каналов в статическом режиме LACP.

\item Настройка некоторых функций статического режима LACP.

\end{itemize}

\subsection{Топология сети}

В лабораторной работе на тему протокола связующего дерева два канала между S1 и
S2 не могли одновременно находиться в состоянии передачи данных. Для
полноценного использования полосы пропускания двух каналов между S1 и S2
необходимо настроить агрегирование каналов Ethernet.

\includegraphics[width=0.5\textwidth]{lb3_61.png} 
\captionof{figure}{Топология сети для лабораторной работы №3}

\subsection{План работы}

\begin{itemize}

\item 1. Настройка агрегирования каналов вручную.

\item 2. Настройка агрегирования каналов в режиме LACP.

\item 3. Изменение параметров для определения активных каналов.

\item 4. Изменение режима балансировки нагрузки.

\end{itemize}

\subsection{Процедура конфигурирования}

\textbf{Шаг 1.} Настройка агрегирования вручную

Создадим Eth-Trunk:

\includegraphics[width=0.5\textwidth]{lb3_62.png} 

Команда interface eth-trunk позволяет перейти в режим существующего Eth-Trunk
или создать Eth-Trunk и перейти в его режим. Цифра 1, используемая в примере,
означает номер порта.

\includegraphics[width=0.5\textwidth]{lb3_63.png} 

Сконфигурируем режим агрегирования каналов для Eth-Trunk:

\includegraphics[width=0.5\textwidth]{lb3_64.png} 

Команда mode позволяет настроить рабочий режим Eth-Trunk, который может быть
LACP или ручная балансировка нагрузки. По умолчанию используется режим ручной
балансировки нагрузки. Таким образом, предыдущую операцию выполнять не
требуется, она приводится только для наглядности.

Добавим порт в Eth-Trunk:

\includegraphics[width=0.5\textwidth]{lb3_65.png}

Для добавления определенного порта в Eth-Trunk можно перейти в режим его
интерфейса и выполнить операцию. Для добавления нескольких портов в Eth-Trunk
можно выполнить команду trunkport в режиме интерфейса Eth-Trunk:

\includegraphics[width=0.5\textwidth]{lb3_66.png}

При добавлении физических портов в Eth-Trunk необходимо учитывать следующие
нюансы:

- Агрегированный канал Eth-Trunk может содержать максимум 8 портов-
участников.

- Eth-Trunk нельзя добавить к другому Eth-Trunk.

- Порт Ethernet можно добавить только к одному Eth-Trunk. Чтобы добавить порт
Ethernet к другому Eth-Trunk, сначала необходимо удалить его из исходного.

- Дистанционные порты, напрямую подключенные к локальным портам-
участникам Eth-Trunk, также должны быть добавлены в Eth-Trunk; в противном
случае два конца не смогут взаимодействовать.

- Такие параметры, как количество физических портов, скорость порта и
дуплексный режим, должны совпадать на обоих концах канала Eth-Trunk.

Выведем на экран статус Eth-Trunk:

\includegraphics[width=0.5\textwidth]{lb3_67.png}

\textbf{Шаг 2.} Настройка агрегирования через  LACP.

Удалим порты-участники из Eth-Trunk:

\includegraphics[width=0.5\textwidth]{lb3_68.png}

\includegraphics[width=0.5\textwidth]{lb3_69.png}

Изменим режим агрегирования на LACP:

\includegraphics[width=0.5\textwidth]{lb3_70.png}

\includegraphics[width=0.5\textwidth]{lb3_71.png}

Команда mode lacp позволяет установить LACP в качестве рабочего режима Eth-
Trunk. Но в моем случае команда lacp-static.

Добавим порт в Eth-Trunk:

\includegraphics[width=0.5\textwidth]{lb3_72.png}

\includegraphics[width=0.5\textwidth]{lb3_73.png}

Выведем на экран статус Eth-Trunk:

\includegraphics[width=0.5\textwidth]{lb3_74.png}

\textbf{Шаг 3.} В обычных условиях в состоянии передачи данных должны находиться только
GigabitEthernet0/0/11 и GigabitEthernet0/0/12, а GigabitEthernet0/0/10 должен
использоваться в качестве резервного порта. Когда количество активных портов
становится меньше 2, Eth-Trunk отключается.

Установим приоритет LACP для S1, чтобы сделать S1 активным устройством:

\includegraphics[width=0.5\textwidth]{lb3_75.png}

Настроим самый высокий приоритет портам GigabitEthernet0/0/11 и
GigabitEthernet0/0/12:

\includegraphics[width=0.5\textwidth]{lb3_76.png}

В режиме LACP пакеты LACPDU (LACP Data Unit) передаются и принимаются обеими
сторонами группы агрегирования каналов.
Сначала выбирается активный инициатор.

Выполняется сравнение полей приоритета системы. По умолчанию используется
значение приоритета 32768 Чем меньше значение, тем выше приоритет. Сторона
с более высоким приоритетом выбирается в качестве активного инициатора
LACP.

При одинаковых приоритетах активным инициатором становится сторона с
меньшим MAC-адресом.

После того, как активный инициатор выбран, устройства на обеих сторонах выбирают
активные порты в соответствии с настройками приоритета порта на активном
инициаторе.

Зададим верхний и нижний пороги активных портов:

\includegraphics[width=0.5\textwidth]{lb3_77.png}

Пропускная способность и статус Eth-Trunk зависят от количества активных портов.

Под пропускной способностью Eth-Trunk подразумевается общая пропускная
способность всех портов-участников в состоянии Up. Для того, чтобы
стабилизировать статус и пропускную способность Eth-Trunk, а также сократить
влияние частых изменений статусов каналов-участников, можно настроить
следующие пороговые значения.

Нижний порог: при сокращении количества активных портов ниже этого порога
Eth-Trunk отключается. Порог определяет минимальную пропускную способность
Eth-Trunk и настраивается с помощью команды least active-linknumber.

Верхний порог: если количество активных портов достигает этого порогового
значения, пропускная способность Eth-Trunk не увеличивается, даже при
увеличении числа каналов в состоянии Up. Верхний порог обеспечивает
доступность сети и настраивается с помощью команды max active-linknumber.

Включим функцию внеочередного занятия линии:

\includegraphics[width=0.5\textwidth]{lb3_78.png}

В режиме LACP при выходе из строя активного канала система выбирает резервный
канал с наивысшим приоритетом, чтобы заменить неисправный. Если включена
функция внеочередного занятия линии, то после восстановления неисправный канал
может снова получить статус активного канала, если он имеет более высокий
приоритет, чем резервный канал. Функцию внеочередного занятия линии можно
включить с помощью команды lacp preempt enable. По умолчанию эта функция
отключена.

Выведем на экран статус текущего Eth-Trunk:

\includegraphics[width=0.5\textwidth]{lb3_79.png}

Отключим GigabitEthernet0/0/12, чтобы смоделировать неисправность канала:

\includegraphics[width=0.5\textwidth]{lb3_80.png}

Отключим GigabitEthernet 0/0/11, чтобы смоделировать неисправность канала:

\includegraphics[width=0.5\textwidth]{lb3_81.png}

\includegraphics[width=0.5\textwidth]{lb3_82.png}

Как можно заметить ethernet-trunk не работает, т.к. в состоянии up находятся меньше, чем 2 порта.

\textbf{Шаг 4.} Изменим режим балансировки нагрузки.

Включим отключенные ранее порты:

\includegraphics[width=0.5\textwidth]{lb3_83.png}

Проверим статус Eth-Trunk 1:

\includegraphics[width=0.5\textwidth]{lb3_84.png}

Функция внеочередного занятия линии включена на Eth-Trunk. Таким образом,
GigabitEthernet0/0/11 и GigabitEthernet0/0/12 становятся активными, потому что имеют
более высокий приоритет, чем GigabitEthernet0/0/10. В результате
GigabitEthernet0/0/10 получает статус Unselect. Кроме того, для обеспечения
стабильности канала время внеочередного занятия линии по умолчанию составляет
30 секунд. Таким образом, внеочередное занятие линии происходит через 30 секунд
после включения портов.

Изменим режим балансировки нагрузки Eth-Trunk на балансировку нагрузки на
основе IP-адреса назначения:

\includegraphics[width=0.5\textwidth]{lb3_85.png}

Чтобы обеспечить правильную балансировку нагрузки между физическими каналами
Eth-Trunk и избежать перегрузки каналов, настройте режим балансировки нагрузки
Eth-Trunk с помощью команды load-balance. Балансировка нагрузки работает только
для исходящего трафика. Поэтому режимы балансировки нагрузки для портов на
разных сторонах виртуального канала могут отличаться.

\subsection{Вопросы.} 

1. Какие требования предъявляются к значениям параметров least active-
linknumber и max active-linknumber?

-least active-linknumber должен быть не меньше 2, а active-linknumber не больше 8ми.

\section{\textit{Лабораторная работа №3.4. Связь между VLAN}}

\subsection{Цели}

Лабораторная работа помогает получить практические навыки по изучению
следующих тем:

\begin{itemize}

\item Использование подинтерфейсов терминирования dot1q для реализации связи
между VLAN

\item Использование интерфейсов VLANIF для реализации связи между VLAN

\item Процесс передачи данных между VLAN

\end{itemize}

\subsection{Топология сети}

Маршрутизаторы R2 и R3 принадлежат к разным VLAN. Для их взаимодействия
необходимы интерфейсы VLANIF и подинтерфейсы терминирования dot1q.

\includegraphics[width=0.5\textwidth]{lb3_87.png} 
\captionof{figure}{Топология сети для лабораторной работы №3}

\subsection{План работы}

\begin{itemize}

\item 1. Настройка подинтерфейсов терминирования dot1q для реализации связи между
VLAN.

\item 2. Настройка интерфейсов VLANIF для реализации связи между VLAN.

\end{itemize}

\subsection{Процедура конфигурирования}

\textbf{Шаг 1.} Настроим основные параметры устройств.

Настроим IP-адреса и шлюзы для R2 и R3:

\includegraphics[width=0.5\textwidth]{lb3_86.png} 

\includegraphics[width=0.5\textwidth]{lb3_88.png} 

\includegraphics[width=0.5\textwidth]{lb3_89.png} 

На S1 назначим R2 и R3 в разные VLAN:

\includegraphics[width=0.5\textwidth]{lb3_90.png} 

\textbf{Шаг 2.} Настроим подинтерфейсы терминирования dot1q для реализации связи между VLAN.

Настроим магистральный порт на S1:

\includegraphics[width=0.5\textwidth]{lb3_91.png} 

Канал между S1 и R1 должен разрешать прохождение пакетов из VLAN 2 и VLAN 3,
потому что R1 должен удалять теги VLAN из пакетов, которыми обмениваются эти
VLAN.

Настроим подинтерфейс терминирования dot1q на маршрутизаторе R1:

\includegraphics[width=0.5\textwidth]{lb3_92.png} 

После создания подинтерфейса осуществляется переход в режим конфигурирования
подинтерфейса. В этом примере цифра 2 указывает номер подинтерфейса.
Рекомендуется, чтобы номер подинтерфейса совпадал с идентификатором VLAN.

\includegraphics[width=0.5\textwidth]{lb3_93.png} 

Команда dot1q termination vid vlan-id позволяет настраивать идентификатор VLAN
для выполнения терминирования Dot1q на подинтерфейсе.

В этом примере, когда GigabitEthernet0/0/1 получает данные с тегами VLAN 2, он
передает данные подинтерфейсу 2 для терминирования VLAN и последующей
обработки. Данные, отправленные с подинтерфейса 2, также помечаются тегами
VLAN 2

\includegraphics[width=0.5\textwidth]{lb3_94.png} 

Подинтерфейсы, выполняющие удаление тегов VLAN, не могут пересылать
широковещательные пакеты и автоматически отбрасывают их при получении. Чтобы
такие подинтерфейсы могли пересылать широковещательные пакеты, необходимо
включить функцию широковещательной передачи ARP с помощью команды arp
broadcast enable. На некоторых устройствах эта функция включена по умолчанию.

\includegraphics[width=0.5\textwidth]{lb3_95.png} 

Проверим связь VLAN:

\includegraphics[width=0.5\textwidth]{lb3_96.png} 

\textbf{Шаг 3.} Настроим интерфейсы VLANIF для реализации связи между VLAN.

Удалим конфигурацию, созданную на предыдущем шаге:

\includegraphics[width=0.5\textwidth]{lb3_97.png} 

\includegraphics[width=0.5\textwidth]{lb3_98.png} 

Создадим интерфейс VLANIF на коммутаторе S1:

\includegraphics[width=0.5\textwidth]{lb3_99.png} 

Снова проверим связь между VLAN:

\includegraphics[width=0.5\textwidth]{lb3_100.png} 

\subsection{Вопросы.} 

1. Какие настройки необходимо выполнить на S1, если маршрутизатору R2
требуется доступ к сети, подключенной к маршрутизатору R1?

-Создать vlan устройств и назначить порты для них, и создать, соответсвенно, Vlanif для них

2. Когда интерфейс VLANIF начинает работать?

-Интерфейс VLANIF перейдет в состояние UP (начнет работать как L3-интерфейс) только при соблюдении двух условий одновременно:

Создан VLAN: В базе данных коммутатора существует соответствующий номер VLAN (команда vlan X).

Активный порт: Хотя бы один физический порт, принадлежащий этому VLAN, находится в состоянии UP:
Это может быть порт в режиме access, которому назначен этот VLAN.
Или порт в режиме trunk, через который этот VLAN разрешен (allow-pass).






